\begin{proposition}[The reduced horizontal case is a minimal elliptic ruled double-curve fibre]\label{cdp:st:C.2.IRR.1}
\Uses{the hypotheses, notation and definitions stated here; \cref{cdp:st:A.5.3}; \cref{cdp:st:A.5.17}; \cref{cdp:st:C4e}; \cref{cdp:st:A.5.15}; \cref{cdp:st:C.2.NC.1}}


Let $f:Y\to C_0$ be a proper holomorphic map from a smooth compact
complex threefold with $H^2(Y,\mathbb Q)=0$, with connected fibres
and smooth general fibre a complex two-torus. Let $mS$ be a full
fibre with irreducible nonnormal reduced support. Let
$h=\eta\pi:Z\to T\to S$ be the minimal resolution of its
normalization, and suppose that $p:Z\to B$ is an irrational
ruling. Write the effective divisor of \ref{cdp:st:A.5.3} as
$D_Z=H+V$, where $V$ is vertical and the horizontal part $H$
is reduced, of total degree two over $B$.

Then $B$ is elliptic, $T=Z$ is a smooth minimal ruled surface
over $B$, and $D_Z=H$ is a disjoint union of smooth elliptic
curves. The reduced conductor image is a disjoint union of
smooth elliptic curves. Its inverse image $H$ maps onto it
by a finite etale covering of total degree two over each
component. The support $S$ has only ordinary double-curve
normal crossings and
\[
                         e(S)=0.                 \tag{C.IR.1}\label{cdp:eq:C.IR.1}
\]
No assumption on the initial normalization singularities or
on the ambient Picard pullback is made.
\end{proposition}
\begin{proof}
\noindent\textbf{Step 1.} The canonical comparison of
\ref{cdp:st:A.5.3} gives
\[
 K_Z+D_Z=h^*\omega_S,\qquad c_1(K_Z+D_Z)=0
 \quad\hbox{in }H^2(Z,\mathbb Q).
 \tag{C.IR.2}\label{cdp:eq:C.IR.2}
\]
If $g(B)\ge2$, that proposition makes the horizontal part
twice a section.  Since the given horizontal part is reduced,
$g(B)=1$.  Contract the vertical ruling blowups to a minimal
ruled surface $\rho:Z\to Z_0$, and write
$D_0=\rho_*D_Z=H_0+V_0$.  Choose a section $C_0$ with
$C_0^2=-n$ and a ruling fibre $F$.  Numerically,
\[
 K_{Z_0}=-2C_0-nF,\qquad H_0=2C_0+bF,\qquad V_0=vF,
 \qquad b+v=n.
\]
Here $v\ge0$.  Adjunction gives $p_a(H_0)=1-v$.
An irreducible $H_0$ has normalization of genus at least one
by Riemann--Hurwitz.  If $H_0$ has two components, they are
sections, and its arithmetic genus is one plus their
intersection number.  Hence $v=0$.  In the irreducible case
$H_0$ is smooth elliptic; in the other case its two smooth
elliptic sections are disjoint.

\noindent\textbf{Step 2.} We justify the numerical comparison
under every blowdown and retain its exact divisor map.
For a point blowup $\sigma:W\to W'$ with exceptional curve
$E$, let $a$ be its coefficient in $D_W$, let
$D'=\sigma_*D_W$, and let $d$ be the multiplicity of $D'$
at the centre.  Then
\[
 K_W+D_W=\sigma^*(K_{W'}+D')+(1+a-d)E.
 \tag{C.IR.4}\label{cdp:eq:C.IR.4}
\]
Intersection with $E$ gives $a=d-1$, so the numerically
trivial canonical-boundary class descends.  This proves
the numerical assertion used in Step~1 without requiring
$h$ to factor through a ruling blowdown.

In the reverse sequence, effectivity forces every centre
to lie on the current divisor: a centre off it would create
coefficient $-1$, which persists on its strict transform.
Starting with the disjoint smooth coefficient-one curves
$H_0$, each permitted centre therefore has $d=1$.  Its
exceptional coefficient is zero, and its new boundary is
the strict transform of $H_0$.  Induction gives
\[
                         D_Z=H,\qquad V=0.
 \tag{C.IR.5}\label{cdp:eq:C.IR.5}
\]
In particular $H$ is a disjoint union of smooth elliptic
curves.  Coefficient-zero ruling exceptional curves have
not yet been discarded.

\noindent\textbf{Step 3.} Let $E$ be a curve exceptional for
the minimal normalization resolution $\pi$ and not contained
in $H$.  Its negative self-intersection and adjunction give
\[
 2p_a(E)-2=E^2-H\cdot E<0.
\]
Thus $p_a(E)=0$, and $E$ is smooth rational.  Minimality
excludes $E^2=-1$.  Since $H\cdot E=2+E^2\ge0$, it follows
that $E^2=-2$ and $H\cdot E=0$.  A positive-coefficient
exceptional curve is one of the mutually disjoint elliptic
components of $H$.  Consequently every connected exceptional
fibre is disjoint from the strict conductor curves.

The normalization conductor ideal is the rank-one reflexive
sheaf $\omega_T\otimes\eta^*\omega_S^{-1}$ embedded in
$\mathcal O_T$ by finite duality.  Any point on its divisorial
support lies on a conductor curve, whose strict transform
must meet the resolution fibre over that point.  The fibres
just described therefore lie outside that support.  There
the conductor ideal equals $\mathcal O_T$, by reflexivity,
and $\eta$ is an isomorphism.  A nontrivial such resolution
fibre would give an isolated normal hypersurface singularity
of $S$.  Proposition~\ref{cdp:st:A.5.17} excludes it in the
actual torus smoothing, with its unit and full multiplicity
retained.  Hence $\pi$ is the identity and $T=Z$ is smooth.

\noindent\textbf{Step 4.} The conductor ideal of $T$ now
defines the reduced smooth divisor $H$.  Its contraction to
$\mathcal O_S$ is radical, since the quotient embeds into
$\mathcal O_H$.  Let $G$ be this reduced conductor image.
At a generic transverse plane curve, conductor exponent one
on a normalization branch forces that branch to be smooth
and to meet exactly one other smooth branch transversely:
the branchwise conductor formula is twice the intrinsic
delta invariant plus the sum of its intersection
multiplicities with the other branches.  This is the
formula used in \ref{cdp:st:C4e}.  Thus $H\to G$ has total
generic degree two over each component.

The sheet involution extends to the smooth compact curves
$H$ over the normalization of $G$; a degree-two map of
smooth curves has its holomorphic deck involution, including
the exchange of two degree-one components.  The involution
preserves the finite fibre over every point of $G$, since
its composites to $G$ agree on a dense open subset and
therefore everywhere.

\noindent\textbf{Step 5.} We prove the local ring comparison
at every point of $G$.  Let $A$ be the reduced hypersurface
ring, let $B$ be its finite normalization ring, and let $I$
be the conductor.  The ring $B$ is a product of regular
two-dimensional local rings.  Define
\[
 A'=\{b\in B:b\bmod I\in(B/I)^\iota\},
\]
where $\iota$ is the actual involution of the preceding
step.  There is an inclusion $A\subset A'\subset B$.
The two first rings agree in codimension one by the
ordinary-node calculation, so $A'/A$ has finite length.
The module $A'$ is finite and torsion-free over $A$.
Local cohomology of $0\to A\to A'\to A'/A\to0$ injects
$A'/A$ into $H^1_{\mathfrak m}(A)$.  The latter vanishes
because a hypersurface surface ring has depth two.
It follows that $A'=A$, and therefore
$A/I=(B/I)^\iota$.

The image curve $G$ is consequently the quotient of a
smooth curve by its involution and is smooth.  The exact
unramified local pullback ring is
\[
 \{(a(x,z),b(y,z)):a(0,z)=b(0,z)\}
       \simeq\mathbb C\{x,y,z\}/(xy).
\]
At a ramification point it is
\[
 \{a(u,t):a(0,t)\in\mathbb C\{t^2\}\}
 =\mathbb C\{u,ut,t^2\}
 \simeq\mathbb C\{x,y,z\}/(x^2-y^2z),
 \tag{C.IR.6}\label{cdp:eq:C.IR.6}
\]
with $(x,y,z)=(ut,u,t^2)$.  Proposition~\ref{cdp:st:A.5.15}
excludes the latter local map, including a unit times the
$m$th power.  Therefore all conductor maps are unramified,
their image components are disjoint smooth elliptic curves,
and $S$ has ordinary double-curve normal crossings.

\noindent\textbf{Step 6.} Apply the normal-crossings
comparison \ref{cdp:st:C.2.NC.1} to the actual full fibre
$mS$.  It gives
\[
 0=e(Y_t)=m\,e(T\setminus H)=m\,e(T),
 \tag{C.IR.9}\label{cdp:eq:C.IR.9}
\]
because $e(H)=0$.  A minimal ruled surface over the elliptic
curve $B$ has Euler number zero, while every point blowup
adds one.  Thus $\rho$ has no blowups and $T=Z$ is minimal.
Finally the actual normalization is an isomorphism off
$H$ and $G$, so $e(S)=e(T)-e(H)+e(G)=0$.
\end{proof}




The elliptic ruled geometry in \ref{cdp:st:C.2.IRR.1} has a direct literature
comparison. Kamenova's Theorem 3.1 in \cite{KamenovaTori} treats
semistable Kahler degenerations with trivial total-space canonical
bundle; its elliptic ruled cases have elliptic double curves arranged
along a path or a cycle. Our finite normalization $T\to S$ specifies
the self-identification of the double curves in the irreducible case.
The proof above establishes this geometry from the actual conductor
and the full-fibre smoothing, retaining the multiplicity and proving
the necessary local pinch exclusion. We have not supplied a Kahler
form or a trivialization of the total-space canonical bundle needed
to apply that theorem directly to our tube. 




